\section{Meaning and voice come first}
\label{sec:principles}

A successful edit changes the prose without changing the writer's claim.
Style is subordinate to factual fidelity. Every name, number, date, quotation,
citation, and technical distinction in the revision must come from the source
or from the writer. When the draft lacks a needed fact, either ask for it or
write the honest, less specific sentence.

Fiction is the exception because invention belongs to the task. In factual,
technical, legal, and reference writing, a vivid unsupported detail remains a
fabrication.

\subsection{Uncertainty is part of the claim}

Preserve the difference between ``demonstrates,'' ``suggests,'' ``is
consistent with,'' and ``does not rule out.'' A shorter sentence is worse when
it turns a conditional result into a universal one. The same principle applies
to opinions. A writer's confidence, hesitation, and mixed response should
survive the edit.

Missing evidence calls for a boundary, not a plausible substitute. If the
available sources do not establish a founding date, say so when the absence
matters or omit the sentence. Do not infer a private life, a motivation, or a
causal history from silence.

\subsection{Voice survives in the choices}

Voice appears in selection as much as style. A writer may keep a technical
qualification that interrupts the rhythm because it is true. Another may use
a dated joke or pause to correct an earlier thought. Do not regularize those
choices merely to make every paragraph sound alike.

When the writer supplies an earlier sample, use it as the strongest evidence
for vocabulary, punctuation, and rhythm. A generic style rule should not
override a consistent choice in that sample unless the choice harms meaning or
violates the publication's requirements.

Voice is often uneven. A short sentence may follow a long one, an aside may
reveal the writer's stake, and a passage may remain ambivalent. Keep that
unevenness when it belongs to the writer.

\subsection{A light edit is often enough}

Read the full draft before touching a sentence. Identify the central point and
the intended reader. Notice a few voice signals worth preserving. Then fix
the actual problems and leave strong sentences alone.

The same restraint applies to structure. An anecdote may create stakes; a
detour may reveal judgment. Keep either when it helps. Reorganize only when
the current order hides priority, repeats a claim, or forces the reader to
reconstruct the logic.

\subsection{Each sentence needs a reason}

Each sentence should give the reader a claim, condition, example, transition,
or moment of voice. This standard is more precise than saying that a sentence
must ``earn its place.'' It identifies the sentence's purpose for the reader.
Throat-clearing, empty qualifiers, and generic encouragement fail this test
because they delay the next useful piece of information.
